Consider a starving forager problem, where an animal (the forager) is searching for food in an environment with a certain distribution of resources. The forager has a limited energy budget and must decide how to allocate its time and energy to maximize its chances of finding food while minimizing the risk of starvation.

We can model this problem as an optimal control problem with the following dynamics and objective:
\begin{align}
    \max J &= \mathbb{E}\left[ \int_{0}^{T} e^{-\rho t} (1 - e^{-\alpha u_{t}}) dt \right] \label{eq:forager:objective}\\
    \text{subject to: } \quad dx_{t} &= (r x_{t} - u_{t})dt + \sigma x_{t} dW_{t} \label{eq:forager:dynamics},
\end{align}
where $x_{t}$ represents the energy level of the forager at time $t$, $u_{t}$ is the foraging effort (control input), $r$ is the natural energy regeneration rate, $\sigma$ is the volatility of the energy level due to environmental factors, $\rho$ is the discount rate, and $\alpha$ controls the saturation of the reward from foraging. The objective function represents the expected discounted reward from foraging over a finite time horizon $T$. The dynamics of the energy level are modeled as a stochastic differential equation (SDE) with a drift term that captures the energy regeneration and depletion due to foraging, and a diffusion term that captures the stochastic fluctuations in energy levels.

\subsection{Baseline/Ground Truth Solution}
